\documentclass[a4paper]{article}
\usepackage{graphicx}
\usepackage{amsmath,amssymb}
\usepackage{hyperref}
\usepackage{minted}
\usepackage{multirow}
\usepackage{float}

\title{Verneinung}
\date{}

\begin{document}
    \maketitle
    
    \subsection{Nicht}
        
        \textbf{Beispiel:}
        \begin{itemize}
            \item Ich habe nicht viel Zeit.
        \end{itemize}
        
        \textbf{Regel:}
        \emph{``nicht`` steht am Ende, wenn man das Verb oder den ganzen Satz verneinen will.}
        
        \begin{itemize}
            \item Ich gehe nicht.
            \item Er kommt nicht.
            \item Wir arbeiten heute nicht.
        \end{itemize}
    
    \subsection{Adjektiv verneinen}
        
        \begin{itemize}
            \item Das ist nicht schön.
            \item Sie ist nicht müde.
            \item Das war nicht einfach.
        \end{itemize}
    
    \subsection{Adverb verneinen}
        
        \begin{itemize}
            \item Er spricht nicht gut Deutsch.
            \item Ich gehe nicht oft ins Kino.
        \end{itemize}
    
    \subsection{Nomen mit Artikel verneinen}
        
        \begin{itemize}
            \item Das ist nicht mein Buch.
            \item Sie ist nicht die Lehrerin.
        \end{itemize}
    
    \subsection{Nomen ohne Artikel verneinen}
        
        \begin{itemize}
            \item Ich habe kein Auto.
        \end{itemize}
    
    
    % \bibliographystyle{plain}
    % \bibliography{/home/donkarlo/Dropbox/repo/shared_research/refs/refs}
\end{document}